\section{وضعیت پردازه، حافظه‌ی مشترک و انتقال پیام}

\subsection*{بخش الف – وضعیت و صف هر پردازه}

در مدل وضعیت‌های درس، پردازه‌ی «در حال اجرا» روی \lr{CPU} است و در هیچ صفی نیست؛ پردازه‌ی «آماده» در صف آماده (\lr{ready queue}) است و پردازه‌ی «در انتظار» در صف انتظارِ همان رویداد یا دستگاهی است که منتظرش است (\lr{device queue})؛ پردازه‌ی «خاتمه‌یافته» دیگر در هیچ‌یک از صف‌های زمان‌بندی نیست. خلاصه‌ی پاسخ‌ها:

\begin{center}\small\atlrows
\begin{tabular}{@{}c>{\raggedleft\arraybackslash}p{3.7cm}>{\raggedleft\arraybackslash}p{8.2cm}@{}}
  \rowcolor{atl-primary}\atlhead{سناریو} & \atlhead{وضعیت} & \atlhead{صف / جایگاه در سیستم‌عامل}\\
  ۱ & در حال اجرا (انتظار فعال) & روی \lr{CPU}؛ در هیچ صفی نیست. با پایان برش زمانی به صف آماده می‌رود، نه صف انتظار\\
  ۲ & در انتظار & صف انتظار نوشتنِ همان لوله (رویداد «فضای خالی در لوله»)\\
  ۳ & خاتمه‌یافته (زامبی) & در هیچ صف زمان‌بندی نیست؛ فقط ورودی‌اش در جدول پردازه‌ها مانده است\\
  ۴ & در حال اجرا (در مد هسته) & روی \lr{CPU}؛ هنوز در هیچ صفی نیست (پس از ذخیره‌ی زمینه به صف آماده می‌رود)\\
\end{tabular}
\end{center}

\textbf{سناریوی ۱ – تولیدکننده با انتظار فعال.} بدنه‌ی حلقه‌ی \code{while} خالی است و پردازه هیچ فراخوانی سیستمیِ مسدودکننده‌ای نمی‌زند؛ پس هیچ‌وقت خودش \lr{CPU} را رها نمی‌کند و از دید سیستم‌عامل کاملاً قابل اجراست. این «انتظار فعال» (\lr{busy waiting}) است: پردازه در حال اجراست (و فقط شرط را بارها محاسبه می‌کند)، نه در حال انتظار. وقتی \lr{CPU} را دارد در هیچ صفی نیست؛ با تمام‌شدن برش زمانی و وقفه‌ی زمان‌سنج به صف آماده می‌رود و نوبت بعدی دوباره اجرا می‌شود. حلقه تا وقتی مصرف‌کننده \code{out} را جلو نبرده باشد تمام نمی‌شود و در این مدت \lr{CPU} هدر می‌رود.

\textbf{سناریوی ۲ – نوشتن مسدودکننده‌ی ۸ کیلوبایت در لوله‌ی بی‌نام.} ظرفیت لوله در لینوکس ۶۴ کیلوبایت است (طبق اسلایدهای درس). چون ۶۰ کیلوبایت هنوز خوانده نشده، فقط ۴ کیلوبایت جای خالی هست و پیام ۸ کیلوبایتی نمی‌گنجد. نوشتن‌های بزرگ‌تر از \code{PIPE\_BUF} (۴ کیلوبایت) اتمیک نیستند: بخشی که جا دارد نوشته می‌شود و پردازه برای بقیه باید صبر کند. چون فراخوانی مسدودکننده است، پردازه خطا نمی‌گیرد و برنمی‌گردد، بلکه می‌خوابد. پس در وضعیت «در انتظار» و در صف انتظار نوشتنِ همان لوله است تا خواننده‌ای فضا باز کند؛ آن‌وقت با رویداد «فضای خالی» بیدار می‌شود و به صف آماده می‌رود. (همین رفتار را روی لینوکس آزموده‌ایم: ۴ کیلوبایت نوشته شد و نویسنده در فراخوانی نوشتن خوابید؛ با خواندن ۸ کیلوبایت از لوله، نوشتن کامل شد.)

\textbf{سناریوی ۳ – فرزندی که پایان یافته ولی والدش منتظرش نشده است.} فرزند \code{exit()} را اجرا کرده و اجرایش تمام شده و حافظه و فایل‌های بازش آزاد شده‌اند، اما چون والد هنوز \code{wait()} نزده و کد خروج را نگرفته، سیستم‌عامل باید بلوک کنترل پردازه را نگه دارد. این وضعیت «خاتمه‌یافته» است و به آن پردازه‌ی زامبی (\lr{zombie}) می‌گویند. چنین پردازه‌ای نه در صف آماده است و نه در صف انتظار، یعنی زمان‌بند هرگز آن را انتخاب نمی‌کند؛ فقط یک ورودی در جدول پردازه‌ها دارد. حلقه‌ی بی‌نهایتِ والد هم چیزی را عوض نمی‌کند، جز اینکه تا وقتی والد زنده است و \code{wait()} نزده، فرزند زامبی می‌ماند (\code{init} هم آن را جمع نمی‌کند، چون والدش هنوز وجود دارد).

\textbf{سناریوی ۴ – ذخیره‌ی ثبات‌ها پس از وقفه‌ی زمان‌سنج.} رفتن به مد هسته وضعیت پردازه را عوض نمی‌کند: کد هسته به‌نام همین پردازه و روی همان \lr{CPU} اجرا می‌شود و \lr{CPU} هنوز در اختیار آن است. ذخیره‌ی ثبات‌ها در \lr{PCB} بخشی از سوئیچ زمینه است و در این لحظه هنوز پردازه به هیچ صفی نرفته است؛ پس وضعیت «در حال اجرا» و بدون صف است. پس از ذخیره‌ی زمینه، اگر برش زمانی تمام شده باشد، زمان‌بند وضعیت را به «آماده» تغییر می‌دهد و \lr{PCB} را به انتهای صف آماده می‌گذارد (گذار «در حال اجرا» به «آماده»).

\subsection*{بخش ب – آرایه‌ی حلقوی و ظرفیت بافر}

در این بخش \lr{$N$} همان \code{BUFFER\_SIZE} است. بافر خالی یعنی \code{in == out} و بافر پر (در پیاده‌سازی کلاسیک) یعنی \code{(in + 1) \% BUFFER\_SIZE == out}.

\textbf{۱) چرا پر کردن هر \lr{$N$} خانه خطا می‌دهد؟} اگر هر \lr{$N$} خانه پر شود، پس از آخرین درج \code{in} یک دور کامل می‌چرخد و دوباره روی \code{out} می‌افتد؛ یعنی حالت «پر» هم \code{in == out} می‌شود و با حالت «خالی» یکی است. الگوریتم فقط از روی \code{in} و \code{out} تصمیم می‌گیرد و دیگر نمی‌تواند این دو را از هم تشخیص دهد.

\textbf{مثال:} \lr{$N=4$} و \code{in = out = 0}. تولیدکننده چهار عنصر می‌گذارد و \code{in} به‌ترتیب ۱، ۲، ۳ و باز ۰ می‌شود. اکنون \code{in == out == 0} در حالی که بافر پر است. مصرف‌کننده شرط \code{in == out} را «خالی» می‌خواند و منتظر می‌ماند و چهار عنصر تولیدشده هرگز مصرف نمی‌شوند. (از طرف دیگر هیچ شرطی نمی‌تواند همان حالت را «پر» بنامد.)

دلیل کلی (اصل لانه‌کبوتری): تفاضل \code{(in - out) \% N} فقط \lr{$N$} مقدار دارد (۰ تا \lr{$N-1$})، ولی تعداد عناصر بافر \lr{$N+1$} حالت دارد (۰ تا \lr{$N$}). پس دو حالت متفاوت ناچار یک‌جور دیده می‌شوند. نگه‌داشتن یک خانه‌ی خالی تعداد حالت‌ها را به \lr{$N$} می‌رساند و هر حالت را یکتا می‌کند.

\textbf{۲) چرا فقط عوض‌کردن شرط تولیدکننده کار نمی‌کند؟} در این راه‌حل (\code{while (in == out)}) شرط «پر بودن» همان شرط «خالی بودن» مصرف‌کننده است. مثال نقض همان حالت آغاز است: مقدار اولیه‌ی \code{in} و \code{out} هر دو صفر است و بافر کاملاً خالی است.
\begin{itemize}
  \item تولیدکننده: \code{in == out} درست است؛ پس فکر می‌کند بافر پر است و در حلقه می‌ماند. هیچ عنصری تولید نمی‌شود.
  \item مصرف‌کننده: \code{in == out} درست است؛ پس بافر را خالی می‌بیند و در حلقه می‌ماند.
\end{itemize}
هیچ‌کدام جلو نمی‌روند و \code{in} و \code{out} هم دیگر تغییر نمی‌کنند، یعنی بن‌بست (که چون هر دو انتظار فعال می‌کنند، \lr{CPU} را هم هدر می‌دهد). اصلاح شرط راه‌حل نیست، چون مشکل در خودِ اطلاعات است: بدون متغیر تازه (مثل شمارنده یا پرچم) هیچ شرطی روی \code{in} و \code{out} نمی‌تواند «پر» و «خالی» را از هم جدا کند.

\subsection*{بخش ج – صندوق پستی مشترک}

\textbf{۱) مدل ارتباط و ابهام.} پردازه‌ها نام یکدیگر را نمی‌دانند و فقط از طریق صندوق پستی \lr{A} حرف می‌زنند؛ پس مدل \textbf{غیرمستقیم} (\lr{indirect}) است (برخلاف مدل مستقیم که در آن \code{send(P, msg)} و \code{receive(Q, msg)} پردازه‌ها را نام می‌برند). ابهام: یک پیام فرستاده شده ولی دو پردازه‌ی \lr{P2} و \lr{P3} هم‌زمان منتظر دریافتند. معلوم نیست پیام به \lr{P2} برسد یا \lr{P3} (یا هر دو، یا هیچ‌کدام)، و فرستنده هم نمی‌داند چه کسی پیام را گرفته؛ نتیجه‌ی برنامه به زمان‌بندی وابسته می‌شود. چهار راه‌حل معماری (مطابق اسلایدهای درس):
\begin{enumerate}
  \item \textbf{هر صندوق حداکثر با دو پردازه مرتبط باشد} (یک فرستنده و یک گیرنده). اگر \lr{P2} و \lr{P3} هر دو گیرنده‌اند، هرکدام صندوق جداگانه می‌گیرند و ابهام ریشه‌کن می‌شود.
  \item \textbf{فقط یک \lr{receive} در هر لحظه.} سیستم‌عامل روی عملیات دریافت هر صندوق انحصار متقابل می‌گذارد؛ اگر \lr{P2} در \code{receive} است، \code{receive} پردازه‌ی \lr{P3} مسدود می‌شود (یا خطا می‌گیرد) تا نوبتش برسد. (پیاده‌سازی ساده: صندوق «مالک» دارد و فقط مالک می‌تواند بگیرد.)
  \item \textbf{سیستم‌عامل گیرنده را انتخاب کند.} هر دو \code{receive} مجازند و سیستم‌عامل یکی را برمی‌گزیند، با قاعده‌ای مثل اولین گیرنده‌ی مسدودشده (\lr{FIFO})، نوبتی (\lr{round-robin}) یا اولویت. پیام فقط به همان یکی تحویل می‌شود و دیگری منتظر پیام بعدی می‌ماند.
  \item \textbf{گیرنده به فرستنده اعلام شود.} سیستم‌عامل شناسه‌ی گیرنده‌ی منتخب را به فرستنده برمی‌گرداند (مثلاً در مقدار بازگشتی یا پیام تأیید) تا فرستنده بداند پیام را چه کسی گرفته است؛ این راه‌حل معمولاً همراه راه‌حل ۳ می‌آید.
\end{enumerate}

\textbf{۲) میعادگاه (\lr{rendezvous}).} میعادگاه یعنی فرستنده و گیرنده باید هم‌زمان «به هم برسند» و هیچ‌کدام بدون دیگری جلو نمی‌رود. این وضعیت وقتی رخ می‌دهد که:
\begin{itemize}
  \item \code{send} مسدودکننده باشد: فرستنده تا دریافت پیام توسط گیرنده منتظر می‌ماند؛
  \item \code{receive} مسدودکننده باشد: گیرنده تا رسیدن پیام منتظر می‌ماند؛
  \item ظرفیت بافر لینک صفر باشد (\lr{zero capacity}): هیچ پیامی در لینک صف نمی‌شود، پس فرستنده فقط با تحویل مستقیم به گیرنده می‌تواند ادامه دهد.
\end{itemize}
اگر ظرفیت لینک \lr{$n>0$} یا نامحدود باشد، \code{send} به‌محض گذاشتن پیام در صف برمی‌گردد و فرستنده به گیرنده وابسته نیست (فقط وقتی صف پر باشد صبر می‌کند)؛ پس میعادگاه کامل رخ نمی‌دهد.

\textbf{چرا تولیدکننده–مصرف‌کننده ساده‌تر می‌شود؟} همگام‌سازی را خودِ \code{send} و \code{receive} (یعنی سیستم‌عامل) انجام می‌دهد:
\begin{codebox}[title=message passing]{C}
/* producer */                       /* consumer */
while (true) {                       while (true) {
    item = produce();                    receive(&item);
    send(item);                          consume(item);
}                                    }
\end{codebox}
\begin{itemize}
  \item پردازه‌ها فضای آدرس جدا دارند و پیام را هسته کپی می‌کند؛ پس متغیر مشترکی نیست و شرایط رقابتی (\lr{race condition}) روی \code{in} و \code{out} یا شمارنده اصلاً پیش نمی‌آید.
  \item بافر دستی، \code{BUFFER\_SIZE}، شاخص‌های \code{in} و \code{out} و محدودیت «\lr{$N-1$} خانه» وجود ندارد.
  \item بررسی پر یا خالی بودن (سرریز و کم‌ریزی) لازم نیست: تولیدکننده نمی‌تواند از مصرف‌کننده جلو بیفتد، چون \code{send} تا دریافت صبر می‌کند، و مصرف‌کننده تا رسیدن پیام خواب است. انتظار هم فعال نیست (پردازه در وضعیت «در انتظار» است و \lr{CPU} هدر نمی‌دهد).
  \item با چند تولیدکننده و مصرف‌کننده هم همین کد درست است، بدون قفل اضافه؛ در حافظه‌ی مشترک برنامه‌نویس باید همه‌ی اینها را خودش پیاده کند و درست هم‌گام کند.
\end{itemize}
(بهای این سادگی، هزینه‌ی فراخوانی سیستمی و کپی‌کردن پیام است.)

\subsection*{بخش و – درخت پردازه‌ها و زمان‌بندی اجرا}

فرض می‌کنیم \lr{P0} در لحظه‌ی ۰ شروع می‌کند و پردازه‌ی ساخته‌شده در تکرار \lr{$i$} حلقه را \lr{C$i$} می‌نامیم. نکته‌های کد:
\begin{itemize}
  \item هر فرزند در پایان کارش \code{exit(0)} می‌زند و تمام می‌شود (\code{exit} داخل بدنه‌ی \code{if} است)، پس فرزندان هرگز به دور بعدی حلقه نمی‌رسند و فقط \lr{P0} حلقه را کامل اجرا می‌کند: در لحظه‌ی ۰ هر سه \code{fork()} را اجرا می‌کند و بعد وارد \code{sleep(3)} می‌شود.
  \item \lr{C1} پس از ۲ ثانیه، \lr{C2} پس از ۴ ثانیه و \lr{C3} پس از ۶ ثانیه از خواب بیدار می‌شود.
  \item \lr{C2} پس از بیدار شدن \code{wait(NULL)} را اجرا می‌کند؛ اما \lr{C2} \textbf{فرزندی ندارد} (\lr{C1} و \lr{C3} خواهرهای او هستند، نه فرزندانش). پس \code{wait} بدون مسدودشدن و با مقدار \lr{$-1$} برمی‌گردد و \lr{C2} فوراً (در لحظه‌ی ۴) \code{exit(0)} می‌زند.
  \item \lr{P0} در لحظه‌ی ۳ از \code{sleep(3)} بیدار می‌شود و \code{wait(NULL)} را \textbf{فقط یک بار} اجرا می‌کند. \lr{C1} از لحظه‌ی ۲ زامبی است، پس \code{wait} فوراً برمی‌گردد و \lr{C1} را جمع می‌کند. (\code{SIGCHLD} در لحظه‌ی ۲ به‌طور پیش‌فرض نادیده گرفته می‌شود و \code{sleep} را قطع نمی‌کند.) سپس \lr{P0} تا لحظه‌ی ۷ می‌خوابد و با \code{return 0} تمام می‌شود.
  \item \lr{P0} دیگر \code{wait} نمی‌زند، پس \lr{C2} (از لحظه‌ی ۴) و \lr{C3} (از لحظه‌ی ۶) تا مرگ \lr{P0} زامبی می‌مانند. در لحظه‌ی ۷ که \lr{P0} می‌میرد، هر دو به \code{init} داده می‌شوند (بازوالدده‌ی) و \code{init} فوراً جمعشان می‌کند.
\end{itemize}

\begin{figure}[ht]\centering
\begin{LTR}
\begin{tikzpicture}[x=1.9cm, y=1cm]
  \node[pnode, minimum width=24mm] (p) at (3,0) {\lr{P0}\\[-1pt]{\scriptsize شروع ۰، پایان ۷}};
  \node[pnode, minimum width=24mm] (c1) at (0.3,-2) {\lr{C1}\\[-1pt]{\scriptsize شروع ۰، پایان ۲}};
  \node[pnode, minimum width=24mm] (c2) at (3,-2) {\lr{C2}\\[-1pt]{\scriptsize شروع ۰، پایان ۴}};
  \node[pnode, minimum width=24mm] (c3) at (5.7,-2) {\lr{C3}\\[-1pt]{\scriptsize شروع ۰، پایان ۶}};
  \draw[pedge] (p) -- (c1); \edgelab{p}{c1}{0.5}{تکرار ۱}
  \draw[pedge] (p) -- (c2); \edgelab{p}{c2}{0.5}{تکرار ۲}
  \draw[pedge] (p) -- (c3); \edgelab{p}{c3}{0.5}{تکرار ۳}
\end{tikzpicture}
\end{LTR}
\caption{درخت پردازه‌ها با زمان شروع و پایان (ثانیه). همه‌ی \code{fork}ها در لحظه‌ی ۰ اجرا می‌شوند.}
\end{figure}

\begin{figure}[ht]\centering
\begin{LTR}
\begin{tikzpicture}[x=1.6cm, y=0.95cm,
  slp/.style={draw=atl-primary, fill=atl-tint, line width=0.5pt},
  zmb/.style={draw=atl-accent, fill=atl-accent!30, line width=0.5pt}]
  \foreach \t in {0,...,7}{
    \draw[gray!45] (\t,-0.55) -- (\t,3.65);
    \node[font=\scriptsize, below] at (\t,-0.55) {\t};
  }
  \node[font=\scriptsize, below] at (8.1,-0.55) {\lr{t (s)}};
  % نام ردیف‌ها
  \node[left] at (-0.15,3) {\lr{P0}};
  \node[left] at (-0.15,2) {\lr{C1}};
  \node[left] at (-0.15,1) {\lr{C2}};
  \node[left] at (-0.15,0) {\lr{C3}};
  % P0
  \path[slp] (0,2.75) rectangle node[font=\scriptsize]{\lr{sleep(3)}} (3,3.25);
  \path[slp] (3,2.75) rectangle node[font=\scriptsize]{\lr{sleep(4)}} (7,3.25);
  \node[tlab, above] at (3,3.3) {\lr{wait()}:\\\lr{C1} جمع می‌شود};
  \node[tlab, above] at (7,3.3) {\lr{return 0}\\\lr{P0} می‌میرد};
  % C1
  \path[slp] (0,1.75) rectangle node[font=\scriptsize]{\lr{sleep(2)}} (2,2.25);
  \path[zmb] (2,1.75) rectangle node[font=\scriptsize]{زامبی} (3,2.25);
  % C2
  \path[slp] (0,0.75) rectangle node[font=\scriptsize]{\lr{sleep(4)}} (4,1.25);
  \path[zmb] (4,0.75) rectangle node[font=\scriptsize]{زامبی تا مرگ \lr{P0}} (7,1.25);
  \node[tlab, above] at (4,1.3) {\lr{wait()} بی‌فرزند:\\فوراً برمی‌گردد، سپس \lr{exit}};
  % C3
  \path[slp] (0,-0.25) rectangle node[font=\scriptsize]{\lr{sleep(6)}} (6,0.25);
  \path[zmb] (6,-0.25) rectangle node[font=\scriptsize]{زامبی} (7,0.25);
  % بازوالدده‌ی در لحظه‌ی ۷
  \node[tlab, right] at (7.05,0.5) {به \lr{init} می‌رسند\\و جمع می‌شوند};
  \draw[->, >=Stealth, thin] (7.05,0.5) -- (7.0,0.95);
  \draw[->, >=Stealth, thin] (7.05,0.5) -- (7.0,0.05);
  % لحظه‌ی ۵
  \draw[atl-accent, very thick, dashed] (5,-0.55) -- (5,3.65);
  \node[font=\small\bfseries, atl-accent, above] at (5,3.65) {\lr{t=5}};
\end{tikzpicture}
\end{LTR}
\caption{وضعیت هر پردازه در طول زمان. رنگ آبی: در خواب؛ نارنجی: زامبی. در لحظه‌ی \lr{t=5} (خط‌چین) فقط \lr{C2} زامبی است.}
\end{figure}

\begin{center}\small\atlrows
\begin{tabular}{@{}c>{\raggedleft\arraybackslash}p{4.4cm}c>{\raggedleft\arraybackslash}p{5.2cm}@{}}
  \rowcolor{atl-primary}\atlhead{پردازه} & \atlhead{ایجادکننده} & \atlhead{پایان} & \atlhead{وضعیت پس از پایان}\\
  \lr{P0} & پردازه‌ی اولیه (آغاز در ۰) & ۷ & —\\
  \lr{C1} & \lr{P0}، تکرار ۱ (در ۰) & ۲ & زامبی از ۲ تا ۳؛ در ۳ با \code{wait} جمع می‌شود\\
  \lr{C2} & \lr{P0}، تکرار ۲ (در ۰) & ۴ & زامبی از ۴ تا ۷؛ در ۷ توسط \code{init} جمع می‌شود\\
  \lr{C3} & \lr{P0}، تکرار ۳ (در ۰) & ۶ & زامبی از ۶ تا ۷؛ در ۷ توسط \code{init} جمع می‌شود\\
\end{tabular}
\end{center}

\textbf{در لحظه‌ی \lr{t=5} ثانیه:} \lr{P0} در \code{sleep(4)} (خواب)، \lr{C3} در \code{sleep(6)} (خواب) و \lr{C2} زامبی است؛ \lr{C1} پیش‌تر (در ۳) جمع شده است.
\begin{enumerate}
  \item \textbf{دقیقاً یک پردازه در وضعیت زامبی است: \lr{C2}.} در لحظه‌ی ۴ مرده، اما والدش \lr{P0} زنده است و دیگر \code{wait} نمی‌زند، پس \lr{PCB} او باقی می‌ماند. \lr{C1} را \lr{P0} در لحظه‌ی ۳ جمع کرده و \lr{C3} هنوز نمرده است.
  \item \textbf{هیچ پردازه‌ی یتیمی نیست (صفر).} بازوالدده‌ی به \code{init} فقط پس از مرگ والد رخ می‌دهد و \lr{P0} تا لحظه‌ی ۷ زنده است؛ پس تا \lr{t=5} هیچ پردازه‌ای به \code{init} سپرده نشده است (چه در حال اجرا و چه خاتمه‌یافته). \lr{C2} و \lr{C3} در لحظه‌ی ۷ یتیم می‌شوند.
\end{enumerate}
